% ============================================================================
%  Rozdział 3. Analiza wymagań i projekt gry
% ============================================================================
\chapter{Analiza wymagań i projekt gry}
\label{rozdz:projekt}

W~rozdziale przedstawiono koncepcję rozgrywki wraz z~zestawieniem elementów występujących
na planszach, wymagania funkcjonalne i~niefunkcjonalne, środowisko pracy i~ograniczenia,
przypadki użycia oraz projekt architektury oprogramowania.

\section{Koncepcja gry i zasady rozgrywki}

Gracz steruje kulą umieszczoną na płaskiej planszy otoczonej ścianami. Kamera śledzi kulę
z~góry i~z~tyłu, pod stałym kątem, dzięki czemu widoczny jest fragment planszy wokół kuli,
a~pełny układ poziomu pokazuje minimapa w~prawym górnym rogu ekranu. Strzał wykonuje się,
dotykając kuli i~przeciągając palec (lub kursor myszy) w~stronę przeciwną do zamierzonego
kierunku ruchu. W~trakcie przeciągania od kuli do palca rysowana jest linia celownicza,
a~po zwolnieniu kula rusza z~mocą proporcjonalną do długości gestu. Kolejny strzał można
wykonać w~dowolnym momencie, także gdy kula jeszcze się toczy.

Na każdym poziomie rozmieszczono monety. Po zebraniu ostatniej z~nich pojawia się meta,
a~wjechanie w~nią kończy poziom wygraną. Kula ma 100~punktów życia. Zderzenie z~kaktusem
odbiera 29, a~z~grzybem 17 punktów; spadek zdrowia poniżej zera oznacza przegraną.
Przeciwnika można też pokonać, trafiając w~jego wrażliwą część, co uruchamia animację
śmierci i~usuwa go z~planszy. Utracone zdrowie można odzyskać, zbierając apteczki,
które przywracają 30~punktów, ale nie więcej niż do wartości początkowej. Apteczka
pozostaje na planszy, jeśli kula przejedzie przez nią z~pełnym zdrowiem. Pozostałe elementy
rozgrywki zestawiono w~tabeli~\ref{tab:elementy}.

\begin{table}[!htb]
  \centering
  \caption{Elementy rozgrywki i~ich parametry; \zrodloWlasne}
  \label{tab:elementy}
  \begin{tabularx}{\textwidth}{>{\raggedright\arraybackslash}p{3.3cm} X >{\raggedright\arraybackslash}p{3.6cm}}
    \toprule
    \textbf{Element} & \textbf{Działanie} & \textbf{Parametry} \\
    \midrule
    Kula & porusza się po strzale, traci moc przez tarcie & moc maks. 25, tarcie 4/s \\
    Ściana & odbija kulę zgodnie z~prawem odbicia & zachowane 70\% mocy \\
    Moneta & dodaje wartość do licznika, odblokowuje metę & wartość 20 \\
    Kaktus & patroluje odcinek pomiędzy dwoma punktami & 29 obrażeń, prędkość 3 \\
    Grzyb & nieruchomy, losowo obrócony & 17 obrażeń \\
    Apteczka & przywraca zdrowie, maks. do 100 & +30 punktów \\
    Pole przyspieszające & wyrzuca kulę wzdłuż strzałek & +20, prędkość maks. 35 \\
    Odbijacz & odpycha kulę od swojego środka & +12, prędkość maks. 30 \\
    Obracająca się belka & obraca się wokół słupka w~losowym kierunku & 90°/s, odepchnięcie 10 \\
    Ruchoma ściana & przesuwa się ruchem harmonicznym & amplituda 2,5, prędkość 2,5 \\
    Wiatr & stale przesuwa kulę w~losowym kierunku & siła 0,3--0,8 \\
    Przycisk \emph{Push} & wypycha stojącą kulę w~losowym kierunku & moc 8 \\
    \bottomrule
  \end{tabularx}
\end{table}

Przygotowano trzy poziomy. Ich zawartość odczytano z~plików scen i~przedstawiono
w~tabeli~\ref{tab:poziomy}. Pierwszy poziom zawiera wyłącznie nieruchomych przeciwników
(grzyby), drugi wyłącznie patrolujące kaktusy, a~trzeci łączy oba typy i~ma dwukrotnie
więcej monet niż pierwszy. Każdy poziom zawiera też jedną metę i~jeden punkt startowy kuli.
Układ pierwszego poziomu w~rzucie z~góry pokazano na rysunku~\ref{rys:poziom1-gora}.

\begin{table}[!htb]
  \centering
  \caption{Liczba elementów na poszczególnych poziomach; \zrodloWlasne}
  \label{tab:poziomy}
  \begin{tabular}{l c c c}
    \toprule
    \textbf{Element} & \textbf{Poziom 1} & \textbf{Poziom 2} & \textbf{Poziom 3} \\
    \midrule
    Wymiary obszaru gry [m] & 76 $\times$ 104 & 104 $\times$ 80 & 108 $\times$ 87 \\
    Segmenty ścian & 29 & 50 & 43 \\
    Monety & 34 & 42 & 82 \\
    Grzyby / kaktusy & 44 / 0 & 0 / 35 & 56 / 33 \\
    Apteczki & 7 & 8 & 8 \\
    Pola przyspieszające & 11 & 12 & 12 \\
    Odbijacze & 11 & 13 & 14 \\
    Obracające się belki & 5 & 6 & 6 \\
    Ruchome ściany & 4 & 5 & 5 \\
    \bottomrule
  \end{tabular}
\end{table}

\begin{figure}[!htb]
  \centering
  \zrzut[0.44\textwidth]{poziom1_widok_z_gory}
  \caption{Rzut z~góry pierwszego poziomu wykonany kamerą ortograficzną. Kolorowe znaczniki
    minimapy oznaczają: niebieski -- kulę, czerwone -- przeciwników, zielone -- apteczki;
    różowe pierścienie to odbijacze, pomarańczowe prostokąty -- pola przyspieszające,
    odcinki z~ciemnym punktem -- obracające się belki; \zrodloWlasne}
  \label{rys:poziom1-gora}
\end{figure}

Na ekranie rozgrywki stale widoczne są licznik monet, czas przejścia poziomu, strzałka
wiatru, minimapa, pasek i~liczba punktów życia, przycisk \emph{Push} oraz panel prędkości
i~kąta obrotu kuli (rysunek~\ref{rys:hud-wstep}).

\section{Wymagania funkcjonalne}

Wymagania funkcjonalne określają działania, które gra musi umożliwiać graczowi, oraz
reakcje systemu na zdarzenia rozgrywki:
\begin{itemize}[leftmargin=1.9cm, itemsep=1pt]
  \item wybór jednego z~trzech poziomów w~menu głównym i~przejście do niego przez ekran
        ładowania,
  \item wykonanie strzału gestem przeciągnięcia rozpoczętym na kuli, z~podglądem kierunku
        w~postaci linii celowniczej,
  \item wyznaczanie mocy strzału na podstawie długości gestu z~ograniczeniem do wartości
        maksymalnej,
  \item odbijanie kuli od ścian i~przeciwników z~utratą części mocy oraz animacją
        zderzenia i~efektem cząsteczkowym przy ścianie,
  \item zbieranie monet i~aktualizacja licznika, a~po zebraniu wszystkich monet
        udostępnienie mety,
  \item zadawanie obrażeń przez przeciwników, odzyskiwanie zdrowia przez apteczki
        i~zakończenie gry przegraną po utracie zdrowia,
  \item działanie przeszkód ruchomych, pól przyspieszających, odbijaczy i~wiatru,
  \item wypchnięcie nieruchomej kuli w~losowym kierunku przyciskiem \emph{Push},
  \item wyświetlanie minimapy, kierunku wiatru, prędkości i~kąta obrotu kuli,
  \item pomiar czasu przejścia poziomu, zapis najlepszego czasu, łącznego czasu gry
        i~największej liczby zebranych monet,
  \item prezentacja statystyk wszystkich poziomów w~oknie dostępnym z~menu głównego,
  \item ponowne rozpoczęcie poziomu, powrót do menu oraz zamknięcie aplikacji.
\end{itemize}

\section{Wymagania niefunkcjonalne}

Wymagania niefunkcjonalne dotyczą technologii, jakości kodu i~parametrów działania
aplikacji:
\begin{itemize}[leftmargin=1.9cm, itemsep=1pt]
  \item implementacja w~silniku Unity 2022.3 LTS w~języku C\#,
  \item działanie na urządzeniach z~systemem Android w~wersji 9 (API~28) lub nowszej,
        z~architekturą ARMv7 lub ARM64, oraz w~edytorze Unity w~systemie Windows,
  \item obsługa wejścia z~ekranu dotykowego i~myszy przez jedną definicję akcji,
  \item płynna rozgrywka; fizyka obliczana ze stałym krokiem 0{,}02~s niezależnie od
        liczby klatek obrazu,
  \item modułowa struktura kodu z~jawnymi zależnościami pomiędzy zestawami kompilacji,
  \item dostarczanie zależności przez kontener Zenject zamiast wyszukiwania obiektów
        w~scenie i~singletonów,
  \item oddzielenie logiki rozgrywki od klas \klasa{MonoBehaviour},
  \item przechowywanie parametrów rozgrywki w~plikach JSON, z~walidacją przed budowaniem,
  \item ograniczenie tworzenia i~niszczenia obiektów w~trakcie rozgrywki,
  \item zwalnianie wszystkich subskrypcji zdarzeń wraz z~obiektami, które je utworzyły,
  \item czytelny interfejs w~orientacji poziomej, skalowany względem rozdzielczości
        referencyjnej 800~$\times$~600.
\end{itemize}

\section{Ograniczenia systemowe i środowisko pracy}

\subsection{Środowisko sprzętowe i programowe}

Prace implementacyjne i~testy w~edytorze prowadzono na komputerze stacjonarnym
z~procesorem AMD Ryzen~5 9600X, kartą graficzną AMD Radeon RX~9070~XT i~32~GB pamięci
operacyjnej, pracującym pod kontrolą systemu Windows~11; dodatkowo grę uruchamiano na
laptopie HP. Wykorzystano edytor Unity
2022.3.60f1 z~modułami Android Build Support, Android SDK \& NDK Tools oraz OpenJDK,
środowisko programistyczne JetBrains Rider oraz system kontroli wersji Git z~repozytorium
w~serwisie GitHub. Biblioteki R3 i~ObservableCollections zainstalowano z~repozytorium
NuGet za pomocą narzędzia NuGetForUnity \cite{nugetforunity}, a~UniTask jako pakiet
menedżera pakietów Unity pobierany z~repozytorium Git. Zenject i~Appodeal dołączono
w~postaci źródłowej w~katalogu \kod{Assets}, a~Firebase jako lokalne pakiety
\kod{.tgz}.

\subsection{Środowisko uruchomieniowe}

Docelową platformą jest system Android. Wersja instalacyjna jest budowana z~zapleczem
IL2CPP, które tłumaczy kod pośredni .NET na C++ i~kompiluje go do kodu natywnego
\cite{unity-il2cpp}, dla architektur ARMv7 i~ARM64. Minimalny poziom API ustawiono na~28,
co odpowiada systemowi Android~9 \cite{android-api}. Interfejs zaprojektowano dla
orientacji poziomej, a~gra nie wymaga połączenia z~siecią.

\subsection{Ograniczenia działania}

Analiza przyjętych rozwiązań pozwoliła wskazać następujące ograniczenia:
\begin{itemize}[leftmargin=1.9cm, itemsep=1pt]
  \item rozpoczęcie strzału wymaga trafienia palcem w~kulę, co na małym ekranie przy
        dużej odległości kamery bywa trudne; łagodzi to przycisk \emph{Push},
  \item wyniki są przechowywane lokalnie w~\klasa{PlayerPrefs}, więc odinstalowanie
        aplikacji je usuwa i~nie są one synchronizowane pomiędzy urządzeniami,
  \item poziomy są zaprojektowane ręcznie i~mają stałą liczbę elementów; losowe są jedynie
        kierunek i~siła wiatru, kierunek obrotu belek oraz faza ruchu ścian,
  \item gra nie zawiera warstwy dźwiękowej,
  \item teksty interfejsu są dostępne wyłącznie w~języku angielskim.
\end{itemize}

\section{Przypadki użycia}

Diagram przypadków użycia przedstawiony na rysunku~\ref{rys:pu} pokazuje funkcje gry
z~perspektywy gracza. Jedynym aktorem jest gracz, a~reakcje gry na zdarzenia rozgrywki
(zebranie monety, otrzymanie obrażeń, ukończenie poziomu) zostały ujęte jako przypadki
rozszerzające przypadek wykonania strzału, ponieważ zachodzą wyłącznie w~jego następstwie.

\begin{figure}[!htb]
  \centering
  \begin{tikzpicture}[diag]
    \aktor{0,0}{Gracz}
    \node[draw, thick, rounded corners, minimum width=11.4cm, minimum height=8.5cm,
          anchor=west, label={[font=\footnotesize\bfseries]above:Gra}] (sys) at (1.3,0) {};
    \node[przypadek, minimum width=3.6cm] (u1) at (4.2, 3.4) {Wybierz poziom};
    \node[przypadek, minimum width=3.6cm] (u2) at (4.2, 2.27) {Przejrzyj statystyki};
    \node[przypadek, minimum width=3.6cm] (u3) at (4.2, 1.13) {Wykonaj strzał};
    \node[przypadek, minimum width=3.6cm] (u4) at (4.2, 0.0) {Wypchnij zablokowaną\\kulę};
    \node[przypadek, minimum width=3.6cm] (u5) at (4.2,-1.13) {Zagraj ponownie};
    \node[przypadek, minimum width=3.6cm] (u6) at (4.2,-2.27) {Wróć do menu};
    \node[przypadek, minimum width=3.6cm] (u7) at (4.2,-3.4) {Wyjdź z~gry};
    \node[przypadek, minimum width=3.3cm] (i1) at (9.7, 3.4) {Wczytaj poziom};
    \node[przypadek, minimum width=3.3cm] (e1) at (9.7, 2.2) {Zbierz monetę};
    \node[przypadek, minimum width=3.3cm] (e2) at (9.7, 1.0) {Otrzymaj obrażenia\\lub ulecz kulę};
    \node[przypadek, minimum width=3.3cm] (e3) at (9.7,-0.2) {Ukończ poziom};
    \foreach \u in {u1,u2,u3,u4,u5,u6,u7} \draw (0.35,0.1) -- (\u.west);
    \draw[strzp] (e1.west) -- node[above, sloped, pos=0.32, font=\scriptsize]{«extend»} (u3.east);
    \draw[strzp] (e2.west) -- node[above, pos=0.32, font=\scriptsize]{«extend»} (u3.east);
    \draw[strzp] (e3.west) -- node[below, sloped, pos=0.32, font=\scriptsize]{«extend»} (u3.east);
    \draw[strzp] (u1.east) -- node[above, font=\scriptsize]{«include»} (i1.west);
    \draw[strzp] (u5.east) -- node[above, pos=0.6, font=\scriptsize]{«include»} (12.15,-1.13) |- (i1.east);
  \end{tikzpicture}
  \caption{Diagram przypadków użycia; \zrodloWlasne}
  \label{rys:pu}
\end{figure}

\subsection{PU Wybierz poziom}

Po uruchomieniu aplikacji i~zainicjalizowaniu usług wyświetlane jest menu główne
z~przyciskami trzech poziomów, statystyk i~wyjścia. Wybór poziomu prowadzi przez ekran
ładowania do poziomu, na którym kula pojawia się w~punkcie startowym, a~licznik czasu
zaczyna odliczanie.

\subsection{PU Wykonaj strzał}

Gracz dotyka kuli i~przeciąga palec, a~gra wyświetla linię od kuli do punktu pod palcem.
Po zwolnieniu gra wyznacza moc i~kierunek strzału; gest rozpoczęty poza kulą jest
ignorowany. W~trakcie ruchu kuli gra realizuje przypadki rozszerzające:
\begin{itemize}[leftmargin=1.9cm, itemsep=1pt]
  \item zebranie monety -- moneta znika z~animacją, licznik monet rośnie, a~po zebraniu
        ostatniej monety pojawia się meta,
  \item otrzymanie obrażeń lub uleczenie -- zmienia się pasek i~liczba punktów życia;
        przy spadku zdrowia poniżej zera kula znika i~wyświetlane jest okno przegranej,
  \item ukończenie poziomu -- wjechanie w~metę zatrzymuje licznik czasu, zapisuje wynik
        i~wyświetla okno końca poziomu z~czasem przejścia.
\end{itemize}

\subsection{PU Wypchnij zablokowaną kulę}

Kula zatrzymana przy ścianie, w~miejscu zasłoniętym dla kamery, może być niemożliwa do
trafienia palcem. Przycisk \emph{Push} jest aktywny wyłącznie wtedy, gdy kula stoi.
Jego naciśnięcie nadaje kuli niewielką prędkość w~losowym kierunku, co pozwala wydostać ją
z~takiego położenia, ale nie daje przewagi w~trakcie normalnej gry.

\subsection{PU Przejrzyj statystyki}

Okno statystyk pokazuje dla każdego poziomu, czy został ukończony, łączny czas gry,
najlepszy czas przejścia oraz największą liczbę monet zebranych w~jednym podejściu.

\subsection{PU Zagraj ponownie, Wróć do menu i Wyjdź z gry}

Okno końca poziomu, wyświetlane zarówno po wygranej, jak i~po przegranej, udostępnia
przycisk ponownego wczytania bieżącego poziomu oraz przycisk powrotu do menu głównego.
Opuszczenie poziomu w~dowolny sposób zapisuje czas spędzony na poziomie. Z~menu głównego
gracz może zamknąć aplikację.

\section{Projekt architektury oprogramowania}
\label{sec:architektura}

Architekturę zaprojektowano tak, aby każda mechanika gry była osobnym modułem, który
otrzymuje zależności z~kontenera, komunikuje się z~innymi modułami przez zdarzenia
i~nie zależy od sposobu, w~jaki zostały utworzone jego współpracujące obiekty. Przyjęto
przy tym zasadę, zgodnie z~którą reguły gry nie zależą od szczegółów technicznych, takich
jak sposób przechowywania konfiguracji czy tworzenia obiektów, lecz to warstwy techniczne
zależą od reguł \cite{martin}. Podział na warstwy przedstawia
rysunek~\ref{rys:architektura}.

\begin{figure}[!htb]
  \centering
  \begin{tikzpicture}[diag]
    \node[modul, minimum width=4.3cm, minimum height=13mm] (json) at (0,0) {\textbf{Konfiguracja}\\ pliki JSON w~Resources\\ \kod{AbstractConfiguration<T>}};
    \node[akcent, minimum width=5.6cm, minimum height=13mm] (infra) at (7.7,0) {\textbf{Infrastruktura}\\ instalatory \kod{MonoInstaller}, bindery,\\ fabryki, \kod{DiContainer}};
    \node[modul, minimum width=4.3cm, minimum height=13mm] (unity) at (0,-2.2) {\textbf{Silnik Unity}\\ PhysX, Input System,\\ URP, Cinemachine};
    \node[akcent, minimum width=5.6cm, minimum height=13mm] (comp) at (7.7,-2.2) {\textbf{Komponenty}\\ klasy \kod{MonoBehaviour}\\ (adaptery zdarzeń silnika)};
    \node[modul, minimum width=4.3cm, minimum height=13mm] (ui) at (0,-4.4) {\textbf{Interfejs (MVVM)}\\ Model, ViewModel, View};
    \node[akcent, minimum width=5.6cm, minimum height=13mm] (logic) at (7.7,-4.4) {\textbf{Logika rozgrywki}\\ klasy \kod{*Logic}, \kod{HealthStorage},\\ \kod{PushPowerConverter}, \dots};
    \node[modul, minimum width=5.6cm, minimum height=11mm] (ext) at (7.7,-6.5) {\textbf{Zapis i~usługi}\\ PlayerPrefs, Firebase, Appodeal};
    \draw[strz] (json) -- node[above, font=\scriptsize]{obiekty danych} (infra);
    \draw[strz] (infra) -- node[right, font=\scriptsize]{wstrzykiwanie \kod{[Inject]}} (comp);
    \draw[strz] (unity) -- node[above, font=\scriptsize]{kolizje, wejście} (comp);
    \draw[strz] (comp) -- node[right, font=\scriptsize]{tworzy, przekazuje zdarzenia} (logic);
    \draw[strz] (logic) -- node[above, font=\scriptsize]{strumienie R3} (ui);
    \draw[strz] (logic) -- node[right, font=\scriptsize]{statystyki poziomu} (ext);
  \end{tikzpicture}
  \caption{Warstwy architektury gry i~kierunki zależności; \zrodloWlasne}
  \label{rys:architektura}
\end{figure}

\subsection{Warstwa infrastruktury}

Warstwa infrastruktury jest korzeniem kompozycji aplikacji. Obiekt \klasa{SceneContext}
każdego poziomu zawiera instalatory poszczególnych obszarów gry (kula, monety, przeciwnicy,
podłoże i~przeszkody, ściany, efekty, linia celownicza, meta, interfejs), które rejestrują
powiązania za pośrednictwem binderów. W~tej warstwie znajdują się też fabryki obiektów.

\subsection{Warstwa konfiguracji}

Parametry liczbowe rozgrywki są zapisane w~18~plikach JSON. Każdej sekcji pliku odpowiada
klasa danych z~polami tylko do odczytu oraz klasa konfiguracji, która ją wczytuje.
Instalatory przekazują obiekty danych do kontenera, więc żadna klasa rozgrywki nie
odczytuje plików samodzielnie.

\subsection{Warstwa komponentów i logiki}

Komponenty \klasa{MonoBehaviour} odbierają zdarzenia silnika (kolizje, kroki fizyki,
wejście) i~przekazują je do klas logiki, które nie dziedziczą po typach silnika,
otrzymują zależności przez konstruktor i~implementują \klasa{IDisposable}, jeśli
subskrybują zdarzenia. Cała reguła gry, np. sposób wyznaczania mocy strzału, znajduje się
więc w~jednej klasie, którą można czytać bez znajomości sceny.

\subsection{Przepływ scen}

Aplikacja składa się ze~sześciu scen (rysunek~\ref{rys:sceny}). Ekran ładowania odczytuje
nazwę poziomu z~obiektu \klasa{LoadLevel} kontekstu projektu, jedynego obiektu
współdzielonego przez wszystkie sceny.

\begin{figure}[!htb]
  \centering
  \begin{tikzpicture}[diag, node distance=5mm and 13mm]
    \node[modul, minimum width=2.2cm] (boot) {\textbf{Bootstrap}\\ Firebase};
    \node[akcent, minimum width=2.6cm, right=of boot] (menu) {\textbf{Main Menu}\\ wybór poziomu,\\ statystyki};
    \node[modul, minimum width=2.2cm, right=of menu] (load) {\textbf{Loading}\\ pasek 1--3 s};
    \node[akcent, minimum width=2.4cm, right=of load] (lvl) {\textbf{Level 1/2/3}\\ rozgrywka};
    \node[modul, minimum width=2.0cm, below=9mm of menu, xshift=-6mm] (quit) {zamknięcie\\ aplikacji};
    \draw[strz] (boot) -- (menu);
    \draw[strz] (menu) -- node[above, font=\scriptsize]{poziom} (load);
    \draw[strz] (load) -- (lvl);
    \draw[strz] ([xshift=-6mm]menu.south) -- node[left, font=\scriptsize]{Quit} ([xshift=-6mm]menu.south |- quit.north);
    \draw[strz] (lvl.south) -- ++(0,-7mm) -| node[pos=0.25, below, font=\scriptsize]{Main Menu} ([xshift=6mm]menu.south);
    \draw[strz] ([xshift=-5mm]lvl.north) -- ++(0,5mm) -| node[pos=0.25, above, font=\scriptsize]{Retry} ([xshift=5mm]lvl.north);
  \end{tikzpicture}
  \caption{Sceny gry i~przejścia pomiędzy nimi; \zrodloWlasne}
  \label{rys:sceny}
\end{figure}
